\section{Resultados algebraicos geométrico--topológicos por caso}
\label{sec:resultados-gt-casos}

Esta sección aplica el atlas de la
\cref{sec:localizacion-geometrico-topologica} a los campos ya fijados en el
reporte. Los resultados presentes son ecuaciones exactas, raíces algebraicas,
simetrías y obstrucciones. La campaña de la \cref{sec:campana-gt} fija cómo
decidir qué partes de las superficies y curvas pertenecen a cada cuenca; la
\cref{sec:resultados-piloto-gt} ejecuta ya su primera ronda B0--B2 para PLL,
MAVPD y Wu. En los casos de Caputo, $F$, $DF\,F$ y
$\mathcal C_{\mathrm{PC}}$ son objetos de arranque en el corte de historias
constantes; no se reinterpretan como velocidad y aceleración globales.

\subsection{Plantilla común de cálculo}

Para cada campo $F=(F_1,\ldots,F_n)^{\mathsf T}$ se ejecutan, en este orden,
las operaciones
\begin{align}
 J(X)&=DF(X), & A(X)&=J(X)F(X),
 \label{casegt:eq:JA}\\
 S_i^v&=\{F_i=0\}, & S_i^a&=\{A_i=0\},
 \label{casegt:eq:surfaces}\\
 S_J&=\{\det J=0\}, &
 \mathcal C&=\{F_iA_j-F_jA_i=0,\ i<j\}.
 \label{casegt:eq:singular-connecting}
\end{align}
Después se resuelve $A=0$, se eliminan las raíces con $F=0$, se comprueba la
región de suavidad y se deduplica por las simetrías. Un PP debe satisfacer
necesariamente $\det J=0$; esta criba evita resolver sistemas polinómicos
innecesarios cuando el Jacobiano es globalmente invertible.

\subsection{Chua PWL entero y fraccionario: superficies por regiones}
\label{subsec:casegt-chua-pwl}

En una región donde la característica es $\phi(x)=mx+d$,
\begin{equation}
 F_1=\alpha[y-(1+m)x-d],\quad
 F_2=x-y+z,\quad F_3=-\beta y-\gamma z.
 \label{casegt:eq:chua-pwl-field}
\end{equation}
El jacobiano regional y su determinante son
\begin{equation}
 J_m=
 \begin{bmatrix}
 -\alpha(1+m)&\alpha&0\\
 1&-1&1\\
 0&-\beta&-\gamma
 \end{bmatrix},
 \qquad
 \det J_m=-\alpha[\beta+(\beta+\gamma)m].
 \label{casegt:eq:chua-pwl-J}
\end{equation}
Multiplicando $J_mF$ se obtiene, sin términos omitidos,
\begin{equation}
 A_1=\alpha[-(1+m)F_1+F_2],\quad
 A_2=F_1-F_2+F_3,\quad
 A_3=-\beta F_2-\gamma F_3.
 \label{casegt:eq:chua-pwl-A}
\end{equation}
Por tanto, las seis superficies regionales son los planos afines
\begin{align}
 S_1^v&:\ y-(1+m)x-d=0,&
 S_1^a&:\ -(1+m)F_1+F_2=0,\nonumber\\
 S_2^v&:\ x-y+z=0,&
 S_2^a&:\ F_1-F_2+F_3=0,\\
 S_3^v&:\ \beta y+\gamma z=0,&
 S_3^a&:\ \beta F_2+\gamma F_3=0.
 \label{casegt:eq:chua-pwl-surfaces}
\end{align}
La curva conectante regional es el conjunto de ceros de
\begin{equation}
 F_1A_2-F_2A_1=0,
 \quad F_1A_3-F_3A_1=0,
 \quad F_2A_3-F_3A_2=0.
 \label{casegt:eq:chua-pwl-CC}
\end{equation}

Con $m_0=-0.1768$, $m_1=-1.1468$ y los parámetros activos,
\begin{equation}
 \det J_{m_0}=-84.035507517056,
 \qquad
 \det J_{m_1}=15.037737580544.
 \label{casegt:eq:chua-pwl-dets}
\end{equation}
Así, $S_J$ es vacío dentro de las tres regiones y no existen PP regionales no
estacionarios. En $x=\pm1$ se evaluaron ambos jacobianos laterales y tampoco se
anula la aceleración lateral. Esto es un control negativo del método PP, pero
no elimina las superficies, las curvas conectantes ni las semillas de función
descriptiva. La simetría central relaciona $X$ con $-X$ y reduce a la mitad las
condiciones iniciales; los cruces $x=\pm1$ se integrarán como eventos exactos.

Para $q=0.9998$, las mismas ecuaciones definen
$\mathcal C_{\mathrm{PC}}$ sólo como atlas de arranque. Las superficies
$F_i=0$ no tienen el teorema de cobertura de una velocidad ordinaria; los
eventos globales se calculan con historia completa.

\subsection{Chua arctangente: superficie singular cerrada y dos escalas}
\label{subsec:casegt-chua-arctan}

Sea
\begin{equation}
 h(x)=a_1x+a_2\arctan(\rho x),
 \qquad
 h'(x)=a_1+\frac{a_2\rho}{1+\rho^2x^2}.
 \label{casegt:eq:arctan-h}
\end{equation}
El campo conserva la forma de \eqref{casegt:eq:chua-pwl-field} sustituyendo
$mx+d$ por $h(x)$. Entonces
\begin{align}
 A_1&=\alpha[F_2-(1+h'(x))F_1],\nonumber\\
 A_2&=F_1-F_2+F_3,\\
 A_3&=-\beta F_2-\gamma F_3,
 \label{casegt:eq:arctan-A}
\end{align}
y
\begin{equation}
 \det J(x)=-\alpha\{\beta+(\beta+\gamma)h'(x)\}.
 \label{casegt:eq:arctan-det}
\end{equation}
La superficie singular se reduce a dos planos $x=\pm x_\star$ cuando
\begin{equation}
 x_\star^2=\frac{1}{\rho^2}
 \left[
 \frac{a_2\rho}{-\beta/(\beta+\gamma)-a_1}-1
 \right]>0.
 \label{casegt:eq:arctan-xstar}
\end{equation}
En $x=x_\star$, defina $p=h'(x_\star)$ y
\begin{equation}
 U=-\frac{\beta x_\star+(\beta+\gamma)h(x_\star)}
 {\gamma+(1+\gamma)p+(\beta+\gamma)/\alpha}.
 \label{casegt:eq:arctan-U}
\end{equation}
Resolver $A=0$ da
\begin{equation}
 y_\star=x_\star+h(x_\star)+\frac{U}{\alpha},
 \qquad
 z_\star=h(x_\star)+(1+p)U+\frac{U}{\alpha}.
 \label{casegt:eq:arctan-yz}
\end{equation}
El segundo punto es el simétrico central. Para Wu, $q=0.99$,
\begin{equation}
 p_\pm=\pm
 \begin{pmatrix}
 0.3369817810952545\\
 0.0815749282686440\\
 -0.2549832342782551
 \end{pmatrix},
 \qquad \lVert F(p_\pm)\rVert=1.39124162716\ldots.
 \label{casegt:eq:wu-pp}
\end{equation}
Para c590, $q=0.9999$,
\begin{equation}
 p_\pm=\pm
 \begin{pmatrix}
 1.213090022614086\\
 -20.18726992409549\\
 -21.54936334065728
 \end{pmatrix},
 \qquad \lVert F(p_\pm)\rVert=545.099830348\ldots.
 \label{casegt:eq:c590-pp}
\end{equation}
Wolfram anula $JF$ en ambos pares a la precisión de trabajo. El primer par fue
integrado causalmente con memoria completa mediante ABM--PECE y EFORK3: ambas
rutas permanecieron acotadas, pero sin destino terminal común ni convergencia
corta al umbral preregistrado, por lo que el resultado llega sólo a EV--TG2
(\cref{sec:resultados-piloto-gt}). El segundo es un control
de escala: una raíz algebraicamente exacta y muy lejana puede ser una semilla
dinámicamente inútil o divergente. La norma de $F$ entra por ello en el filtro
antes de integrar.

\subsection{Kalman--Fitts: sin PP, pero con sección crítica natural}
\label{subsec:casegt-kalman}

Expanda
\begin{equation}
 p(s)=(s^2+2\beta s+m_1^2+\beta^2)
       (s^2+2\beta s+m_2^2+\beta^2)
 =s^4+a_3s^3+a_2s^2+a_1s+a_0.
 \label{casegt:eq:kalman-poly}
\end{equation}
Con $(m_1,m_2,\beta)=(0.9,1.1,0.03)$,
\begin{equation}
 (a_0,a_1,a_2,a_3)
 =(0.98191881,0.121308,2.0254,0.12).
 \label{casegt:eq:kalman-coefficients}
\end{equation}
En coordenadas compañeras,
\begin{equation}
 F=(x_2,x_3,x_4,g)^{\mathsf T},
 \quad
 g=-a_0x_1-a_1x_2-a_2x_3-a_3x_4
 -\tanh(x_3/\varepsilon).
 \label{casegt:eq:kalman-field}
\end{equation}
Si
\begin{equation}
 \kappa(x_3)=a_2+\varepsilon^{-1}
 \operatorname{sech}^2(x_3/\varepsilon),
 \label{casegt:eq:kalman-kappa}
\end{equation}
la multiplicación por el jacobiano produce
\begin{equation}
 A=(x_3,x_4,g,H)^{\mathsf T},
 \quad
 H=-a_0x_2-a_1x_3-\kappa(x_3)x_4-a_3g.
 \label{casegt:eq:kalman-A}
\end{equation}
Así,
\begin{align}
 S_1^v&:\ x_2=0,& S_1^a&:\ x_3=0,\nonumber\\
 S_2^v&:\ x_3=0,& S_2^a&:\ x_4=0,\nonumber\\
 S_3^v&:\ x_4=0,& S_3^a&:\ g=0,\\
 S_4^v&:\ g=0,& S_4^a&:\ H=0.
 \label{casegt:eq:kalman-surfaces}
\end{align}
La sección $x_3=0$ usada en el mapa de conmutación del expediente es a la vez
$S_2^v$ y $S_1^a$: el método existente ya estaba utilizando una superficie
geométrica privilegiada.

Además,
\begin{equation}
 \det J=a_0=(\beta^2+m_1^2)(\beta^2+m_2^2)
 =0.98191881>0.
 \label{casegt:eq:kalman-det}
\end{equation}
No existe $S_J$ ni PP no estacionario. La curva conectante sí puede existir;
sus seis ecuaciones son
\begin{gather}
 x_2x_4-x_3^2=0,
 \quad x_2g-x_3x_4=0,
 \quad x_2H-x_3g=0,\nonumber\\
 x_3g-x_4^2=0,
 \quad x_3H-x_4g=0,
 \quad x_4H-g^2=0.
 \label{casegt:eq:kalman-CC}
\end{gather}
Se parametrizarán sus componentes reales dentro del dominio acotado y se
compararán contra la semilla del mapa de conmutación. El campo es impar, así
que sólo se conserva un representante de cada par.

\subsection{MAVPD: superficie singular, PP exactos y fallo informativo}
\label{subsec:casegt-mavpd}

Para
\begin{equation}
 F_1=\delta(\gamma u+v-u^3),\quad
 F_2=u-\xi v-w,\quad F_3=\rho v,
 \label{casegt:eq:mavpd-field}
\end{equation}
se obtiene
\begin{align}
 A_1&=\delta[(\gamma-3u^2)F_1+F_2],\nonumber\\
 A_2&=F_1-\xi F_2-F_3,\\
 A_3&=\rho F_2,
 \label{casegt:eq:mavpd-A}
\end{align}
y las superficies
\begin{align}
 S_1^v&:\ \gamma u+v-u^3=0,&
 S_1^a&:\ (\gamma-3u^2)F_1+F_2=0,\nonumber\\
 S_2^v&:\ u-\xi v-w=0,&
 S_2^a&:\ F_1-\xi F_2-F_3=0,\nonumber\\
 S_3^v&:\ v=0,& S_3^a&:\ F_2=0.
 \label{casegt:eq:mavpd-surfaces}
\end{align}
Como
\begin{equation}
 \det J=\delta\rho(\gamma-3u^2),
 \qquad
 S_J=\left\{u=\pm\sqrt{\gamma/3}\right\},
 \label{casegt:eq:mavpd-det}
\end{equation}
resolver $A=0$ fuera de los equilibrios da exactamente
\begin{equation}
 u_\star=\pm\sqrt{\frac\gamma3},
 \quad
 v_\star=\frac{2\delta\gamma u_\star}{3(\rho-\delta)},
 \quad
 w_\star=u_\star-\xi v_\star.
 \label{casegt:eq:mavpd-PP}
\end{equation}
Los tres menores $F_iA_j-F_jA_i$ definen la curva conectante y contienen estos
dos puntos. Para $\delta=100$, $\rho=200$, $\gamma=0.1$, $\xi=3.1$,
\begin{equation}
 p_\pm=\pm(0.1825741858351,0.0121716123890,0.1448421874291).
 \label{casegt:eq:mavpd-PP-numeric}
\end{equation}
El piloto B0--B2 lleva estos PP a los ciclos internos y las semillas DF1 a la
rama exterior. La persistencia de horizonte no alcanza el umbral: la deriva
B1--B2 es $26.24\%$ para los PP y $10.69\%$ para DF1. Además, el bracket entre
ambas rutas encuentra un tercer destino, de modo que no define todavía una
separatriz binaria. El máximo conservador es EV--TG2 y no EV--TG4
(\cref{sec:resultados-piloto-gt}). Es un resultado negativo útil: confirma que
$\operatorname{PP}\subset\mathcal C$ no cubre todas las cuencas. La campaña
ampliada deberá muestrear las superficies completas, las ramas de $\mathcal C$
y el borde entre ciclos. La simetría $X\mapsto-X$ reduce el banco a un hemisferio;
los parámetros $\xi=3.5$ y el candidato caótico en $\xi=2.85$ funcionarán como
control negativo y extensión, respectivamente.

\subsection{PLL: curvas en el cilindro y banco KCC natural}
\label{subsec:casegt-pll}

Con $T=\tau_1+\tau_2$ y el campo de \cref{int:eq:pll}, escriba
$F=(F_x,F_\theta)^{\mathsf T}$. Entonces
\begin{align}
 A_x&=\frac1T\left[-F_x+\frac{\tau_1}{2}
 \cos\theta\,F_\theta\right],\nonumber\\
 A_\theta&=-\frac LT\left[F_x+\frac{\tau_2}{2}
 \cos\theta\,F_\theta\right].
 \label{casegt:eq:pll-A}
\end{align}
Las superficies críticas, que ahora son curvas sobre
$\R\times S^1$, quedan
\begin{align}
 S_x^v&:\ x=\frac{\tau_1}{2}\sin\theta,&
 S_x^a&:\ -F_x+\frac{\tau_1}{2}\cos\theta F_\theta=0,\nonumber\\
 S_\theta^v&:\ x=\frac{T\omega_\Delta}{L}
 -\frac{\tau_2}{2}\sin\theta,&
 S_\theta^a&:\ F_x+\frac{\tau_2}{2}\cos\theta F_\theta=0.
 \label{casegt:eq:pll-surfaces}
\end{align}
El determinante
\begin{equation}
 \det J=\frac{L\cos\theta}{2T}
 \label{casegt:eq:pll-det}
\end{equation}
anula en los dos círculos
$\theta=\pi/2,3\pi/2\pmod{2\pi}$. Las intersecciones no estacionarias de
$A_x=A_\theta=0$ son
\begin{equation}
 p_1=(\tau_1/2,\pi/2),
 \qquad p_2=(-\tau_1/2,3\pi/2),
 \label{casegt:eq:pll-PP}
\end{equation}
si $\omega_\Delta\mp L/2\ne0$. La curva conectante es
\begin{equation}
 -LF_x^2+
 \left(1-\frac{L\tau_2}{2}\cos\theta\right)F_xF_\theta
 -\frac{\tau_1}{2}\cos\theta\,F_\theta^2=0.
 \label{casegt:eq:pll-CC}
\end{equation}
No hay simetría central para $\omega_\Delta\ne0$; sólo la transformación de
cubierta $\theta\mapsto\theta+2\pi k$. En B0--B2, $p_+$ converge de manera
robusta al foco y alcanza EV--TG3, mientras $p_-$ queda como transitorio acotado
sin destino y en EV--TG2. El bracket propuesto se rechaza porque sus extremos
siguen sin resolver; no existe trayectoria de borde ni EV--TG4. El ciclo silla
y el ciclo estable ya localizados continúan siendo controles para una ventana
posterior más larga, el retorno intervalar y el invariante KCC de
\eqref{gt:eq:pll-kcc-P}; véase la \cref{sec:resultados-piloto-gt}.

\subsection{Lorenz generalizado fraccionario: dos semillas FPP--A}
\label{subsec:casegt-lorenz}

Con $(\sigma,r,a)=(17/5,34/5,-1/2)$,
\begin{equation}
 F=\left(
 -\frac{17}{5}x+\frac{17}{5}y+\frac12yz,
 \frac{34}{5}x-y-xz,
 -z+xy
 \right)^{\mathsf T}.
 \label{casegt:eq:lorenz-field}
\end{equation}
Las tres superficies $S_i^v$ son los ceros de los tres componentes anteriores.
La multiplicación exacta $A=DF\,F$ da
\begin{align}
 A_1={}&\frac{867}{25}x-\frac{374}{25}y+\frac{1}{2}xy^2
 -\frac{27}{10}yz-\frac12xz^2,\nonumber\\
 A_2={}&-\frac{748}{25}x+\frac{603}{25}y-x^2y
 +\frac{27}{5}xz-\frac12yz^2,\\
 A_3={}&\frac{34}{5}x^2-\frac{27}{5}xy+\frac{17}{5}y^2
 +z-x^2z+\frac12y^2z.
 \label{casegt:eq:lorenz-A}
\end{align}
Por tanto $S_i^a$ son las tres superficies algebraicas
$A_i=0$, y $\mathcal C_{\mathrm{PC}}$ se obtiene de los tres menores
$F_iA_j-F_jA_i$. El determinante general es
\begin{align}
 \det J={}&-ay^2+2axyz+az^2-axyr-azr-\sigma
 -\sigma x^2-\sigma xy-\sigma z+\sigma r.
 \label{casegt:eq:lorenz-det}
\end{align}

Wolfram resolvió $A=0$: tres raíces son los equilibrios y las dos restantes
son
\begin{equation}
 p_1=(-0.745966861002653,12.66180451381419,-7.442305306160726),
 \quad
 p_2=(-p_{1x},-p_{1y},p_{1z}).
 \label{casegt:eq:lorenz-PP}
\end{equation}
El residuo máximo de $A$ es menor que $10^{-45}$ y
$\lVert F(p_i)\rVert=23.42210895339\ldots$. La simetría
$(x,y,z)\mapsto(-x,-y,z)$ permite integrar un representante por perfil de
memoria. En $q=0.995$ estas raíces son semillas FPP--A; todavía no son una
reproducción dinámica local. Antes de buscar una frontera entre $E_+$, $E_-$ y
el candidato se exige un conjunto absorbente y control de no explosión.

\subsection{Rabinovich--Fabrikant: ocho semillas en cuatro clases}
\label{subsec:casegt-rf}

Con
\begin{align}
 F_1&=y(z-1+x^2)+ax,\nonumber\\
 F_2&=x(3z+1-x^2)+ay,\\
 F_3&=-2z(b+xy),
 \label{casegt:eq:rf-field}
\end{align}
el jacobiano es
\begin{equation}
 J=\begin{bmatrix}
 2xy+a&z-1+x^2&y\\
 3z+1-3x^2&a&3x\\
 -2yz&-2xz&-2(b+xy)
 \end{bmatrix}.
 \label{casegt:eq:rf-J}
\end{equation}
En forma factorizada, las aceleraciones son
\begin{align}
 A_1&=(2xy+a)F_1+(z-1+x^2)F_2+yF_3,\nonumber\\
 A_2&=(3z+1-3x^2)F_1+aF_2+3xF_3,\\
 A_3&=-2yzF_1-2xzF_2-2(b+xy)F_3.
 \label{casegt:eq:rf-A}
\end{align}
Estas expresiones definen $S_i^a$ sin expandir innecesariamente el polinomio.
Con $(a,b)=(0.1,0.2876)$, $A=0$ tiene 13 raíces reales: cinco equilibrios y ocho
PP. Cuatro representantes son
\begin{align}
 p_1&=(-1.04880884817015,1.14415510709471,0),\nonumber\\
 p_2&=(0.948683298050514,0.843274042711568,0),\nonumber\\
 p_3&=(0.217907153070161,0.326271267123146,1.202380810758825),\\
 p_4&=(-0.197739086435437,0.226277627330655,0.822145202940591).
 \label{casegt:eq:rf-PP}
\end{align}
Los otros cuatro se obtienen mediante
$(x,y,z)\mapsto(-x,-y,z)$. Sus normas de campo son, respectivamente,
\begin{equation}
 0.21950\ldots,
 \quad0.17951\ldots,
 \quad1.34455\ldots,
 \quad0.76891\ldots.
 \label{casegt:eq:rf-PP-speeds}
\end{equation}
Wolfram verifica residuos menores que $10^{-40}$ y pertenencia a $S_J$. Para
$q=0.998$, la campaña usa cuatro representantes, sus perfiles perturbados y un
destino explícito de escape. Por los términos cúbicos, un bloque acotado debe
preceder cualquier afirmación de recurrencia o índice.

\subsection{Muñoz--Pacheco: cuatro raíces exactas sin equilibrios}
\label{subsec:casegt-munoz}

El campo publicado es
\begin{equation}
 F=(yz+x(y-a),\ 1-|x|,\ -xy-z)^{\mathsf T}.
 \label{casegt:eq:munoz-field}
\end{equation}
Fije una región $x\ne0$ y $s=\operatorname{sign}x$. Entonces
\begin{equation}
 J_s=\begin{bmatrix}
 y-a&z+x&y\\
 -s&0&0\\
 -y&-x&-1
 \end{bmatrix},
 \label{casegt:eq:munoz-J}
\end{equation}
y
\begin{align}
 A_1&=(y-a)F_1+(z+x)F_2+yF_3,\nonumber\\
 A_2&=-sF_1,\\
 A_3&=-yF_1-xF_2-F_3.
 \label{casegt:eq:munoz-A}
\end{align}
Además,
\begin{equation}
 \det J_s=s(-x+xy-z).
 \label{casegt:eq:munoz-det}
\end{equation}
De $A_2=0$ se obtiene $F_1=0$. Como no hay equilibrios, $F_2\ne0$ en una
raíz PP; $A_3=0$ implica $F_3=-xF_2$. Al sustituir en $A_1=0$ resulta
$z=x(y-1)$, y de $F_1=0$ se deduce $y^2=a$. Si
\begin{equation}
 u_\varepsilon=2(1-\varepsilon\sqrt a),
 \qquad \varepsilon\in\{-1,1\},
 \label{casegt:eq:munoz-u}
\end{equation}
las cuatro raíces exactas son
\begin{equation}
 (x,y,z)=
 (s u_\varepsilon,\ \varepsilon\sqrt a,\
 -s u_\varepsilon^2/2),
 \qquad s\in\{-1,1\}.
 \label{casegt:eq:munoz-PP}
\end{equation}
Para $a=0.35$ se obtienen
\begin{align}
 &(\pm0.816784043380077,\ 0.591607978309962,
 \ \mp0.333568086760154),\nonumber\\
 &(\pm3.183215956619923,\ -0.591607978309962,
 \ \mp5.066431913239846).
 \label{casegt:eq:munoz-PP-numeric}
\end{align}
Sus normas de campo son $0.2365640964\ldots$ y $7.2845066702\ldots$. Todas
están separadas de $x=0$, por lo que el cálculo regional es válido. La simetría
es $(x,y,z)\mapsto(-x,y,-z)$. Para $q=0.97$ y $0.996$ estas raíces son semillas
Picard--Caputo. Como no existen equilibrios, la definición usual de ocultedad
por vecindades de equilibrios es vacía; aquí se estudiarán localización,
coexistencia, aislamiento y dependencia de la historia, no se presentará la
ausencia de equilibrios como prueba de caos.

\subsection{Resumen de lo obtenido antes de integrar}

\begin{longtable}{@{}L{3.1cm}L{3.0cm}L{3.0cm}L{4.2cm}@{}}
\caption{Resultado algebraico actual y pregunta dinámica siguiente.}
\label{tab:casegt-summary}\\
\toprule
Caso & Superficie singular/PP & Simetría útil & Pregunta que queda abierta\\
\midrule
\endfirsthead
\toprule
Caso & Superficie singular/PP & Simetría útil & Pregunta que queda abierta\\
\midrule
\endhead
Chua PWL & $S_J=\varnothing$ regional; sin PP & Central, con intercambio de
regiones & ¿Qué componentes de $S_i^v,S_i^a,\mathcal C$ caen en la cuenca
oculta?\\
Chua arctan Wu & Dos PP moderados & Central & FPP--A acotados, sin destino común;
EV--TG2 y prueba de convergencia pendiente.\\
Chua c590 & Dos PP con $\lVert F\rVert\approx545$ & Central & ¿Escapan bajo
todo refinamiento o existe una ruta de borde?\\
Kalman--Fitts & $\det J>0$; sin PP & Central & ¿La curva conectante recupera el
ciclo hallado por conmutación?\\
MAVPD & Dos PP exactos & Central & PP interiores y DF exterior; tercer destino
en el bracket, máximo EV--TG2.\\
PLL & Dos PP sobre círculos singulares & Cubierta angular & $p_+$ alcanza
EV--TG3; $p_-$ y el borde quedan sin resolver.\\
Lorenz $q=0.995$ & Dos FPP--A & $(-x,-y,z)$ & ¿Existe región absorbente y borde
funcional entre $E_\pm$ y el candidato?\\
RF $q=0.998$ & Ocho FPP--A, cuatro clases & $(-x,-y,z)$ & ¿Qué semillas son
acotadas, recurrentes y distintas de los cinco equilibrios?\\
Muñoz $q=0.97/0.996$ & Cuatro FPP--A exactos & $(-x,y,-z)$ & ¿Cómo se separan
las cuencas funcionales de los estados caótico y periódico publicados?\\
\bottomrule
\end{longtable}

\ifHAFOdigitalcontent
\subsection{Trazado de la validación algebraica}

La reconstrucción independiente se conserva en
\codepath{wolfram/topologia_geometria_hidden_validation.wl} y su salida legible
en \codepath{wolfram/topologia_geometria_hidden_validation.txt}. Su contraparte
integrada al repositorio es
\codepath{version_2/validation/wolfram/cases/geometric_topological_engine.wl}.
La corrida contiene 50 comprobaciones aprobadas y ninguna fallida. Los
resultados por caso no se obtuvieron leyendo las nubes dinámicas: los campos,
jacobianos, determinantes, simetrías y raíces se reconstruyeron desde las
ecuaciones. El alcance de esta evidencia es algebraico. Los artefactos de la
primera campaña dinámica están trazados por separado en la
\cref{sec:resultados-piloto-gt}.
\fi
